\section{Related Work}
\label{sec:related}

\subsection{Hate-speech detection with transformers}

The application of pretrained transformer models to hate-speech detection has become the dominant paradigm since BERT~\citep{devlin2019bert}. \citet{liu2019roberta} introduced RoBERTa with extended pretraining and dynamic masking; \citet{he2023debertav3} proposed DeBERTa-v3, whose disentangled attention separately encodes content and position and shows consistent NLU gains. Domain-specific pretraining has also been explored: \citet{caselli2021hatebert} retrained BERT on banned Reddit communities to produce HateBERT, outperforming general-purpose BERT on OffComBR and AbusEval. A recurring failure mode is \emph{identity-term shortcut learning}: \citet{vidgen2020directions} and \citet{dixon2018measuring} document that the mere presence of words like ``Muslim,'' ``gay,'' or ``Jews'' can trigger positive predictions independent of context. \citet{dixon2018measuring} introduced the FPED and FNED metrics we report in \S\ref{sec:fairness} to quantify this disparity.

\subsection{Antisemitism detection in NLP}

Antisemitism detection has received comparatively little dedicated NLP attention. \citet{jikeli2023antisemitic} created the GoldStandard dataset of IHRA-annotated English tweets, which we use in its 2024 version~\citep{jikeli2024goldstandard}. \citet{patel2025evaluating} conducted a comprehensive LLM evaluation (GPT-4, Claude, Llama families) under zero-shot, few-shot, and Guided Chain-of-Thought prompting, finding that even the best-prompted LLMs lag behind fine-tuned classifiers on overall macro-F\textsubscript{1} but excel in qualitative reasoning. We replicate their Guided CoT strategy across four LLM families. \citet{becker2025decoding} provide a broader perspective on the linguistic complexity of antisemitic discourse, including coded language and emerging multimodal challenges. \citet{liu2024detecting} apply LoRA fine-tuning to antisemitism detection across BERT, DistilBERT, RoBERTa, and Llama-2; they study accuracy but not counterfactual fairness or calibration.

\subsection{Counterfactual Logit Pairing and counterfactual fairness for text}
\label{sec:related:clp}

The closest prior art to our CCI v2 loss is the Counterfactual Logit Pairing (CLP) family. \citet{garg2019counterfactual} introduced CLP at AIES, framing it as ``counterfactual token fairness'' enforced through a paired-logit penalty in L\textsubscript{p} space; they explicitly note (\S2) that CLP is an instance of individual fairness via Lipschitz constraint in the sense of~\citet{dwork2012fairness}. \citet{davani2021improving} (WOAH; not 2023, as occasionally cited) adapted CLP to hate-speech detection on the Gab Hate Corpus, adding a language-model-likelihood filter to reject implausible swaps; their $\lambda$ hyperparameter governs the strength of the pair-invariance penalty. \citet{mou2024fairness} apply counterfactual logit pairing in a fairness-aware hate-speech framework with explicit uncertainty estimation. \citet{kaushik2020learning} establish the broader counterfactual-augmentation principle on sentiment data, showing that human-edited counterfactuals reduce shortcut reliance and improve distribution-shift robustness; \citet{sen2023people} document a \emph{negative} result --- LLM-generated counterfactuals underperform human-generated ones on sexism and hate-speech detection --- which motivates our use of a deterministic, heuristic-filter-based swap engine rather than an LLM generator. CCI v2 generalises this family along two specific axes: divergence in \emph{probability} space rather than L\textsubscript{p} in logit space (\S\ref{sec:cci}), and a per-group temperature head that we report does not contribute beyond fixed-$T$ on CFR\textsubscript{hard}. We position our methodological contribution as an empirical audit of probability-space counterfactual regularization rather than a fundamentally new fairness algorithm.

\subsection{Concept erasure and representation debiasing}
\label{sec:related:erasure}

An alternative family of approaches operates at the representation rather than at the loss level. \citet{ravfogel2020null} introduced Iterative Null-space Projection (INLP), which iteratively trains linear classifiers to predict a protected attribute from representations and projects them onto the null-space, removing linearly-extractable identity information. \citet{belrose2023leace} proposed LEACE (Least-squares Concept Erasure), a closed-form linear concept erasure with theoretical no-leakage guarantees; LEACE has become the canonical post-hoc representation debiasing tool. Both methods require access to inner-layer representations and are tokenizer-aware (the projection is fit on the specific embedding geometry). CCI v2 differs structurally: it operates at the \emph{text} level (the symmetric vocabulary is text, the tokenizer of each backbone tokenises it natively), which makes the intervention tokenizer-agnostic. We do not include LEACE/INLP in our main results table because our pipeline trains classifiers from scratch rather than starting from a pretrained-representation baseline; a head-to-head LEACE-post-hoc comparison on cached [CLS] representations is a natural follow-up.

\subsection{Calibration, multicalibration, and fairness incompatibilities}
\label{sec:related:calibration}

Modern neural classifiers are over-confident~\citep{guo2017calibration}, and Temperature Scaling is the canonical post-hoc fix. \citet{mukhoti2020calibrating} establish that focal loss yields better-calibrated networks than cross-entropy at fixed accuracy; \citet{tao2023pcs} document that ``early stopping fails to calibrate networks,'' which underpins our bimodal-ECE observation (\S\ref{sec:bimodality}). \citet{hebertjohnson2018multicalibration} introduced multicalibration at ICML, requiring calibration within every subgroup induced by a computationally bounded family of indicator functions. \citet{petersen2021post} study post-processing for individual fairness specifically in large-scale NLP/BERT settings, providing a methodological contrast to our training-time approach.

The fairness-incompatibility literature is older than our specific finding. \citet{pleiss2017fairness} prove that calibration is incompatible with equalized false-positive and false-negative rates across groups. Our \S\ref{sec:multicalibration_tension} finding is narrower and complementary: empirically, post-hoc multicalibration --- a tool designed to enforce subgroup-conditional calibration --- can substantially \emph{degrade} the L\textsubscript{1}-norm soft counterfactual flip rate (an individual-fairness metric) on a text classifier (+96\% increase). \citet{hansen2024multicalibration} provide one of the few empirical multicalibration studies in NLP, examining subgroup ECE and accuracy across tabular, image, and language datasets but \emph{not} counterfactual-pair invariance. To our reading, our finding is the first quantitative report of this specific multicalibration vs.\ counterfactual-invariance trade-off on a text classifier.

\subsection{LLM prompting and statistical-significance protocols}

Large language models exhibit strong zero-shot and few-shot classification ability~\citep{brown2020language}, but performance is sensitive to prompt phrasing, example selection, and output parsing. Chain-of-Thought prompting~\citep{wei2022chain} encourages step-by-step reasoning before the final classification; the faithfulness of those reasoning traces is itself contested~\citep{turpin2023language}. 4-bit NormalFloat quantization~\citep{dettmers2022gptint} makes 7--8B-parameter models tractable on a single consumer GPU. We benchmark four LLM families --- Mistral-7B-Instruct~\citep{jiang2023mistral}, Qwen-2.5-7B-Instruct~\citep{qwen25}, Llama-3.1-8B-Instruct~\citep{llama31}, and Llama-3.3-70B (via the Groq inference API) --- under zero-shot, few-shot, and Guided CoT prompts. Finally, multi-comparison significance is essential when evaluating along multiple connected axes. We follow \citet{dror2018hitchhikers} in using paired-bootstrap p-values over the seed dimension, organised into two pre-specified hypothesis families with Holm-Bonferroni correction per family and an exact sign-test backup for the small-$n$ regime (\S\ref{sec:sig_protocol}).
